\documentclass[10pt,a4paper]{amsart}
\usepackage[margin=19mm]{geometry}
\usepackage{amsmath,amssymb,mathtools}
\usepackage[scr=rsfs,cal=euler]{mathalfa}
\usepackage{xfrac}
\usepackage[unicode,hidelinks,bookmarks=false]{hyperref}
\newcommand{\ie}{i.e.\ }
\newcommand{\cf}{cf.\ }
\newcommand{\resp}{respectively\ }
\newcommand{\Subsection}{\subsection}
\providecommand{\nobreakdash}{\nobreak\hbox{-}}
\setlength{\emergencystretch}{2em}
\multlinegap0pt
\usepackage{bm}
\usepackage{bbm}
\usepackage{dsfont}
\usepackage{xspace}
\usepackage{cancel}
\usepackage{mathtools}
\usepackage{amsthm}
\makeatletter
\@ifundefined{theorem}{\newtheorem{theorem}{Theorem}}{}
\@ifundefined{lemma}{\newtheorem{lemma}[theorem]{Lemma}}{}
\makeatother
\makeatletter
\expandafter\ifx\csname AMS\endcsname\relax
\newenvironment{AMS}{\vspace{.2cm}\small \noindent \begin{quote}\textbf{MSC codes.}}{\end{quote}}
\fi
\expandafter\ifx\csname assumption\endcsname\relax
\newenvironment{assumption}[1][]{%
  \refstepcounter{assumptioncounter}%
  \trivlist
  \item[\hskip\labelsep\hskip\parindent\scshape Assumption~\theassumptioncounter.]%
  \normalfont\ignorespaces
}{\endtrivlist}
\fi
\renewcommand{\geq}{\geqslant}
\expandafter\ifx\csname keywords\endcsname\relax
\newenvironment{keywords}{\small \noindent\begin{quote}\textbf{Keywords.}}{\end{quote}}
\fi
\renewcommand{\leq}{\leqslant}
\newcommand{\mc}[1]{\mathcal{#1}}
\newcommand{\mescin}{\mu}
\newcommand{\norm}[1]{\left\lVert#1\right\rVert}
\newcommand{\slangle}[1]{\left\langle #1 \right\rangle}
\makeatother
\title[On the diffusive-mean field limit of a kinetic weakly intera]{On the diffusive-mean field limit of a kinetic weakly interacting particle system}
\author{Rapha\"el Gastaldello, Grigorios A. Pavliotis, Gabriel Stoltz, Urbain Vaes}
\date{Source: arXiv:2607.15896v1 (CC BY 4.0)}
\begin{document}
\maketitle

\noindent\emph{Source and licence.} On the diffusive-mean field limit of a kinetic weakly interacting particle system. Version arXiv:2607.15896v1 (CC BY 4.0), distributed under the Creative Commons Attribution 4.0 International licence, as recorded in the arXiv licence metadata for this exact version and shown on the abstract page https://arxiv.org/abs/2607.15896v1. Full rights notice: licenses/arxiv-2607.15896-CC-BY-4.0.txt.\par\smallskip

\noindent\textit{Editorial scope.} Counted unit: the Lemma whose source label is \texttt{lem:born\_inf\_norm\_H1} in the article (source0.tex lines 1907--1914, source label \texttt{lem:born\_inf\_norm\_H1}), delivered together with the article's own complete original proof of it (source0.tex lines 1915--1991, one \texttt{proof} environment of 77 lines). No mathematics is written, repaired or re-derived: statement and proof are the article's own text line for line. Required context retained uncounted: eq:U\_N\_norm\_sym\_mat (lines 1894--1897). Every definition, display and label of those lines is retained unchanged. Omitted from this package and not redistributed: 1--1893, 1898--1906, 1992--3268. Nothing inside a retained excerpt is omitted or altered: every retained body is delivered exactly as the article's lines are written, so no omission receipt of any kind accompanies this package. Citations: the retained excerpts carry no citation command at all, so no bibliography entry of the article is transcribed and no reference is redistributed. Reference adaptation: none was needed and none was made; every reference inside the delivered unit resolves inside it, so no unresolved reference or citation is produced. Printed slips kept literal: no printed word, symbol, hypothesis or number of the counted statement or of its proof was altered. Layout and class: the article's own class is \texttt{siamart251216} with packages including authblk, silence, fontenc, inputenc, csquotes, babel, paralist, enumitem, xcolor, graphicx, bbm, amsmath, amssymb, amsfonts; the wrapper is the lane's pinned \texttt{amsart} editorial wrapper at 10pt on A4, which restates the theorem-like environments the retained text needs and adds the article's own macro definitions that the retained text needs (\texttt{\textbackslash geq}, \texttt{\textbackslash leq}, \texttt{\textbackslash mc}, \texttt{\textbackslash mescin}, \texttt{\textbackslash norm}, \texttt{\textbackslash slangle}), with extra wrapper packages (bm, bbm, dsfont, xspace, cancel, mathtools). The delivered numbering is the unit's own. Assets: no retained span contains a figure environment, a graphics-inclusion command, a PDF-page inclusion, a table or a TikZ picture; no image file is copied into this package, the package declares no asset dependency and no image file is redistributed with it. Licence: version arXiv:2607.15896v1 of this article is distributed under the Creative Commons Attribution 4.0 International licence, as recorded in the arXiv licence metadata for this exact version and shown on the abstract page https://arxiv.org/abs/2607.15896v1. A full rights notice is delivered at licenses/arxiv-2607.15896-CC-BY-4.0.txt. Characters: every retained excerpt is pure ASCII, so the delivered engine, XeLaTeX with its default Latin Modern text font, reports no missing character.
    \begin{align}
        \label{eq:U_N_norm_sym_mat}
        -(K_W+K_V) I_{Nd} \leq \nabla^2 U_N \leq (K_W+K_V) I_{Nd}.
    \end{align}

\begin{lemma}
    \label{lem:born_inf_norm_H1}
    There exists~$a_*>0$ such that, for any~$a \in (0,a_*)$, there is~$\kappa_a >0$ for which
    \begin{align*}
        \forall N,\qquad \forall u\in H^1(\mescin_N),\qquad\slangle{-\mc{L}_N u ,u}_a \geq \kappa_a \left(\norm{\nabla_p u}^2+\norm{\nabla_q u}^2\right),
    \end{align*}
    with~$\kappa_a$ only depending on~$\gamma$, $\beta$,~$d$,~$V$ and~$W$.
\end{lemma}

\begin{proof}
  Note that
\begin{align}
    \label{eq:minoration}
\slangle{-\mc{L}_N u ,u}_a = \frac{\gamma}{\beta}\norm{\nabla_p u}^2
+ a\slangle{-(\nabla_p - \nabla_q)\mc{L}_N u ,(\nabla_p - \nabla_q)u}\!.
\end{align}
Straightforward computations give
\begin{align}
    \label{eq:commutators}
    \nabla_q \mc{L}_N = -\nabla^2U_N\nabla_p + \mc{L}_N\nabla_q, \\
    \nabla_p \mc{L}_N = \nabla_q - \gamma\nabla_p + \mc{L}_N\nabla_p,
\end{align}
where~$\mc{L}_N\nabla_p$ is to be understood as the operator~$\mc{L}_N$ applied to each coordinate of~$\nabla_q$.
Inserting~\eqref{eq:commutators} into~\eqref{eq:minoration} gives
\begin{align*}
    \slangle{-\mc{L}_N u ,u}_a = &\frac{\gamma}{\beta}\norm{\nabla_p u}^2 + a\slangle{-\mc{L}_N(\nabla_p - \nabla_q) u ,(\nabla_p - \nabla_q)u} \\
    &+ a \slangle{(\gamma -\nabla^2U_N) \nabla_pu,(\nabla_p - \nabla_q)u} + a\slangle{-\nabla_qu,(\nabla_p - \nabla_q)u}\!.
\end{align*}
Since
\[
\slangle{-\mc{L}_N(\nabla_p - \nabla_q) u ,(\nabla_p - \nabla_q)u} = \frac{\gamma}{\beta}\norm{\nabla_p(\nabla_p - \nabla_q)u}^2 \geq 0,
\]
the following lower bound holds:
\begin{align*}
    \slangle{-\mc{L}_N u ,u}_a &\geq \frac{\gamma}{\beta}\norm{\nabla_p u}^2
    + a \slangle{(\gamma -\nabla^2U_N) \nabla_p u,(\nabla_p - \nabla_q)u} \\
    &\quad + a\slangle{-\nabla_q u,(\nabla_p - \nabla_q)u}.
\end{align*}
Next, by~\eqref{eq:U_N_norm_sym_mat}, and denoting by~$K = K_W + K_V$,
\begin{align}
    \label{eq:hypo_vil1}
a\slangle{(\gamma -\nabla^2U_N) \nabla_p u,(\nabla_p - \nabla_q) u} \geq a(\gamma - K)\norm{\nabla_p u}^2 - a(\gamma + K)\norm{\nabla_p u}\norm{\nabla_q u}.
\end{align}
and
\begin{align}
    \label{eq:hypo_vil2}
a\slangle{-\nabla_q u,(\nabla_p - \nabla_q) u} \geq a \norm{\nabla_q  u}^2 - a \norm{\nabla_p u}\norm{\nabla_q u}.
\end{align}
% By the Cauchy--Schwarz inequality applied to the second terms on the right-hand sides of~\eqref{eq:hypo_vil1} and~\eqref{eq:hypo_vil2}, we have
% \[
% - a \slangle{(\gamma +\nabla^2U_N) \nabla_p u, \nabla_q u} - a \slangle{ \nabla_p u, \nabla_q u}\geq -a(\gamma + K +1)\norm{\nabla_p  u}\norm{\nabla_q  u}.
% \]
Finally,
\[
\slangle{-\mc{L}_N  u , u}_a \geq \left(\frac{\gamma}{\beta} + a(\gamma - K)\right)\norm{\nabla_p  u}^2 + a \norm{\nabla_q u}^2 -a(\gamma + K +1)\norm{\nabla_p u}\norm{\nabla_q u}.
\]
With~$X = (\norm{\nabla_p u},\norm{\nabla_q u})^\top$, the latter inequality can be reformulated as
\[
\slangle{-\mc{L}_N u ,u}_a \geq X^\top\mc{M}_aX,
\]
with
\[
\mc{M}_a :=\begin{pmatrix}
            \displaystyle\frac{\gamma}{\beta} + a(\gamma - K) & \displaystyle-\frac{a}{2}(\gamma + K + 1) \\
            \displaystyle-\frac{a}{2}(\gamma + K + 1) & a
            \end{pmatrix}.
\]
The next step of the proof is to show the existence of an~$a >0$ such that the smallest eigenvalue of~$\mc{M}_a$ is positive, which is equivalent to requiring
$$
\text{Tr}(\mc{M}_a) = \frac{\gamma}{\beta} + a(\gamma - K +1) >0,
$$
and
$$
\text{Det}(\mc{M}_a) = a\left(\frac{\gamma}{\beta} + a(\gamma - K)\right)-\frac{a^2}{4}(\gamma + K + 1)^2>0.
$$
The first inequality is satisfied when~$a \in \left(0,\displaystyle\frac{\gamma}{\beta(K - \gamma -1)}\right)$ if~$\gamma - K +1 <0$, or when~$a>0$ if~$\gamma - K +1 \geq 0$. 
The second inequality is equivalent to
\[
\frac{\gamma}{\beta}a^{-1}+\gamma - K > \frac14(\gamma + K + 1)^2.
\]
Satisfying these two inequalities is of course possible for~$a>0$ sufficiently small since~$\gamma/\beta >0$. In conclusion, there exists $a_*>0$ such that, for any~$a \in (0,a_*)$, there is~$\kappa_a >0$ so that
\[
\slangle{-\mc{L}_N u ,u}_a \geq \kappa_a \left(\norm{\nabla_p u}^2+\norm{\nabla_q u}^2\right),
\]
which concludes the proof.
\end{proof}

\end{document}
